The rapid proliferation of Large Language Models (LLMs) such as Llama 3, Qwen3, and Mixtral has catalyzed a paradigm shift in automated knowledge synthesis. However, current research tools remain fragmented: a researcher must manually query multiple databases, cross-reference findings, synthesize narratives, and generate visualizations---a process consuming 40--60 hours per comprehensive literature review. This project presents the \textbf{AI Research Assistant}, a production-grade, multi-agent autonomous research system that leverages Groq's ultra-low-latency LLM inference API (processing at 500+ tokens/second on Llama-3.3-70B) to execute end-to-end research pipelines without human intervention.

The system orchestrates five specialized AI agents operating in a coordinated pipeline: (1)~\textbf{Research Agent}---a parallel tournament-based system that spawns three concurrent LLM instances (historical, state-of-the-art, and emerging research perspectives) with varying temperatures ($T \in \{0.2, 0.4, 0.6\}$), evaluates each output via an LLM-as-a-judge scoring mechanism, and selects the highest-quality research brief (targeting scores $\geq$~90/100); (2)~\textbf{Methodology Agent}---generates PhD-caliber technical methodology proposals with specific architectures, loss functions, and quantitative targets for each identified research gap; (3)~\textbf{Visualizer Agent}---produces publication-quality system architecture diagrams, workflow visualizations, and methodology flowcharts using Pillow-based programmatic rendering with 8 randomized visual themes; (4)~\textbf{Invoice/Report Agent}---synthesizes all collected research into a 6000+ word PhD-grade whitepaper with inline citations, literature survey tables, and comparative analysis; and (5)~\textbf{AutoResearch Agent}---implements Karpathy's autonomous experimentation paradigm for iterative code optimization with automated evaluation loops.

The frontend is built on Streamlit with a WebSocket-aware polling architecture that maintains stable connections during long-running multi-agent pipelines (up to 10+ minutes). The system supports configurable research domains (computer science, physics, biology, economics), adjustable temporal scope (1--20 year lookback), and six distinct system architecture categories. Key technical innovations include robust JSON extraction with truncation repair, multi-model failover chains (Qwen3-32B $\rightarrow$ Llama-3.1-8B $\rightarrow$ Llama-3.3-70B), and a self-correcting report generation loop inspired by Karpathy's autoresearch methodology. Output artifacts include downloadable Markdown reports, PDF documents, and high-resolution PNG architecture diagrams rendered at 200 DPI across 3200$\times$2000+ pixel canvases.
